% Folder 136, batch 16 (archivists' pages 301-320).
% Pass: Opus 5.5 (claude-opus-5-5), 2026-10-03 — first pass, unchecked against the pages by a human.
%
% Page 318 is a blue cover in his hand (« Gribouillis vrac », « Au Fil des Jours… »):
% no \page, its words become the section heading before pages 319-320.
\documentclass[11pt,a4paper]{article}
\input{../preamble/grothendieck}

\folder{136}
\batch{16}
\pages{301}{320}
\foldertitle{Complexe de De Rham à puissance divisée [conférence de 1976 à l’IHÉS] : notes manuscrites (s.d.).}
\dating{[à partir de 1975-1976]}
\watermark{Édition de démonstration}

\begin{document}

\note{Les pages 301 à 317 forment une suite paginée par l'auteur de 1 à 17, qui commence avec ce lot ; elle ne semble pas reprendre un argument du lot précédent. La p. 317 s'achève sur la fin d'un paragraphe ; on ne sait pas si la suite continue au-delà}

\page{301}
\note{folio de l'auteur (1) ; numéro au crayon des archivistes : 303. Une nouvelle suite paginée par l'auteur commence ici}
$\Delta$ = cat. des simplexes $\Delta^{n}$ ($n \geqslant 0$) \qquad
$\Delta^{n} = [0, n] \subset \mathbb{N}$ avec rel. d'ordre totale

$\widehat{\Delta} = \mathcal{S}s = \underline{\mathrm{Hom}}(\Delta^{\circ}, (\mathrm{Ens})) = \mathrm{Ens}_{*}$

Plus généralement, si $\mathcal{A}$ est une catégorie, on pose
\[
  \mathcal{A}_{*} = \underline{\mathrm{Hom}}(\Delta^{\circ}, \mathrm{Ens}), \qquad
  \mathcal{A}^{*} = \underline{\mathrm{Hom}}(\Delta, \mathrm{Ens})
\]
\note{« $\mathrm{Ens}$ » tel quel dans ces deux formules, là où l'on attend $\mathcal{A}$}
donc
\[
  (\mathcal{A}^{\circ})^{*} = (\mathcal{A}_{*})^{\circ}, \qquad
  \mathcal{A}^{*} \simeq ((\mathcal{A}^{\circ})_{*})^{\circ}
\]
où l'exposant $\circ$ désigne la \add{\ill{}} cat. opposée.

\marginal{La notation $\widehat{\Delta}$ peut être utilisée de préf. quand on veut garder à l'esprit la nature \emph{algébrique} très particulière de ce topos, la notation $\mathcal{S}s$ quand \struck{on} veut l'oublier au profit de sa signification topologique.}

Les objets de $\mathcal{A}_{*}$ (resp. $\mathcal{A}^{*}$) sont considérés comme des
structures \emph{algébriques} de type simple (sur une infinité d'objets de base) dans
$\mathcal{A}$, les « objets semi-simpliciaux » resp. « semi-cosimpliciaux ». On écrit
$\mathrm{Ens}_{*}$ quand on a en vue cet aspect, et $\widehat{\Delta}$ \add{ou
$\mathcal{S}s$} (quand on a plutôt le p\supplied{oin}t de vue « objet d'un topos »
ou « faisceau sur un topos »). Tous les développements qui suivent ont pour objet
d'étudier ces objets, (l'équivalent combinatoire des espaces topologiques) et leurs
« types d'homotopie », via des invariants qui leur sont associés, qui en dépendent
de façon covariante ou contravariante. Nous avons ici en vue une étude
systématique, dans l'esprit de l'algèbre universelle, de ces invariants, (qui sont
donc des foncteurs sur $\widehat{\Delta} = \mathrm{Ens}_{*}$) \add{de leurs
structures et des} \struck{et des} opérations qui peuvent être définies sur elles,
permettant d'en déduire certaines \uncertain{complexes} à partir d'autres plus
simples
\note{« semi-simpliciaux », « semi-cosimpliciaux » : il écrit ½simpliciaux, ½cosimpliciaux, au sens ancien de « simpliciaux »}

\page{302}
\note{folio de l'auteur (2) ; numéro au crayon : 304}
Un objet de $\mathcal{A}_{*}$ sera généralement noté par un symbole de la forme
$K_{*}$, où $K_{*}$ désigne la famille des
\[
  K_{*}(\Delta^{n}) = K_{[n]} \qquad K_{*} = (K_{[n]})_{n \geqslant 0}
\]
\add{\uncertain{par abus de notation}} \add{ou mieux $\Delta^{n} \longmapsto K_{[n]}$}
avec les opérations semi-simpliciales entre elles. On fera attention qu'on écrit
$K_{[n]}$ et non $K_{n}$, pour des raisons qui vont apparaître \uncertain{limpidement}
lorsque $\mathcal{A}$ est additive …).
De même un objet de $\mathcal{A}^{*}$ sera noté
\[
  K^{*} = (n \longmapsto K^{[n]}) \overset{\text{abus.}}{=} (K^{[n]})_{n}
\]

Nous aurons à travailler avec la situation, où on a deux catégories $\mathcal{A}$ et
$\mathcal{B}$, et une équivalence (notée \struck{\ill{}} $M \mapsto \check{M}$)
\[
  v : \mathcal{A}^{\circ} \xrightarrow{\ \approx\ } \mathcal{B} \quad \text{d'où} \quad
  \mathcal{B}^{\circ} \xrightarrow{\ \approx\ } \mathcal{A}
\]
[le plus souvent $\mathcal{A} = \mathcal{B}$ et $(\check{M}^{\vee})^{\vee} \simeq M$
\add{isom. fonct.}, avec compatibilités habituelles d'une autodualité…] on notera
alors \add{souvent} par la même lettre un objet de $\mathcal{A}_{*}$ (ou
$\mathcal{A}^{*}$) et l'objet de $\mathcal{A}^{*}$ (resp. $\mathcal{A}_{*}$) qui lui
correspond par application de $v$, mais en indiquant la variance par la position
du signe $*$
\[
  K^{*} = (K_{*})^{\vee} \qquad K^{[n]} = K_{[n]}{}^{\vee}
\]
\[
  K_{*} = (K^{*})^{\vee} \qquad K_{[n]} = K^{[n]\vee}
\]
On s'intéressera surtout au cas où ($k$ étant un anneau fixé) \struck{$\mathcal{A}$ \ill{}}
\[
  \text{\struck{cat. $\mathcal{A}$}} \quad
  \begin{aligned}
    \mathcal{A} &= k\text{-modules à g. proj. de t.f.} \\
    \mathcal{B} &= k\text{-modules à droite proj. de t.f.}
  \end{aligned}
\]
\note{après « $\mathcal{A} =$ », un mot biffé et illisible}
(si $k$ est commutatif, on a $\mathcal{A} = \mathcal{B}$ et on a une autodualité
sur $\mathcal{A}$).

On a le foncteur canonique pl. fid.
\[
  (*) \qquad \Delta \hookrightarrow \widehat{\Delta} = \mathcal{S}s
\]
l'image par ce foncteur de $\Delta^{n}$ se note \add{ou mieux $\Delta^{[n]}{}_{*}$
(puisqu'il est covariant en $n$)} $\Delta^{n}{}_{*}$. On a donc
\[
  \Delta^{n}{}_{*} \;(\text{ou } \Delta^{[n]}{}_{*}) = (m \longmapsto \Delta^{n}{}_{[m]}
  = \mathrm{Hom}_{\mathrm{Ens}}(\Delta^{m}, \Delta^{n}))
\]
et pour $X_{*} \in \mathrm{Ob}\,\widehat{\Delta}$,
\[
  \mathrm{Hom}(\Delta^{n}{}_{*}, X_{*}) \simeq X_{*}(\Delta^{n}) \simeq X_{[n]}
\]
en particulier
\[
  \mathrm{Hom}(\Delta^{n}{}_{*}, \Delta^{m}{}_{*}) \simeq \Delta^{m}{}_{*}(\Delta^{n}{}_{*})
  \simeq \mathrm{Hom}_{\Delta}(\Delta^{n}, \Delta^{m})
  \quad (= \mathrm{Hom}_{\mathrm{ord}}(\Delta^{n}, \Delta^{m}))
\]
\note{l'indice de $\mathrm{Hom}_{\mathrm{Ens}}$ et celui de $\mathrm{Hom}_{\Delta}$ sont lus avec doute ; $\mathrm{Hom}_{\mathrm{ord}}$ est écrit sous la ligne}
Pour $n$ variable, $\Delta^{n}{}_{*} \in \widehat{\Delta}$ dépend de façon
\emph{covariante} de $n$, \struck{i.e.} donc le foncteur $n \mapsto \Delta^{n}{}_{*}$
(qui n'est autre que $(*)$) est un objet canonique de $\widehat{\Delta}^{*}$, qu'on
appelle \emph{objet cosimplicial universel} et note $\Delta^{*}$ (ou mieux
$\Delta^{*}{}_{*}$)
\[
  \Delta^{*}{}_{*} \in \widehat{\Delta}^{*} = (\mathrm{Ens}_{*})^{*} \quad \text{donc } \forall n \geqslant 0
\]
\[
  \Delta^{[n]}{}_{*} \in \widehat{\Delta} \qquad
  \Delta^{[n]}{}_{*} = \Delta^{n}{}_{*} = m \longmapsto \Delta^{n}{}_{[m]}
  = \mathrm{Hom}_{\Delta}(\Delta^{m}, \Delta^{n})
\]
\struck{\ill{}}
[Une parenthèse de notations : considérons \uncertain{des} catégories de la forme
$(\mathcal{A}_{*})^{*}$, un objet de ladite est donc noté
$K^{*} = (K^{[n]})_{n}$, où les $K^{[n]}$ sont des $\in \mathcal{A}_{*}$, donc
s'écrivent $K^{[n]}{}_{*}$,
\[
  K^{[n]}{}_{*} = (K^{[n]}{}_{[m]})_{m \geqslant 0}
\]
Ainsi, on est amené à noter $K^{*}{}_{*}$ les objets de
$(\mathcal{A}_{*})^{*} \simeq \underline{\mathrm{Hom}}(\Delta \times \Delta^{\circ}, \mathcal{A})$
\uncertain{on}

\page{303}
\note{folio de l'auteur (3) ; numéro au crayon : 305}
\[
  K^{[m]}{}_{[n]} = K(\Delta^{m}, \Delta^{n}) \in \mathcal{A}
\]
De même, on note
\[
  \begin{aligned}
    &K_{*}{}^{*} \ \text{les objets de } (\mathcal{A}^{*})_{*} \simeq \underline{\mathrm{Hom}}(\Delta^{\circ} \times \Delta, \mathcal{A}) \\
    &K^{**} \ \text{les objets de } (\mathcal{A}^{*})^{*} \simeq \underline{\mathrm{Hom}}(\Delta^{\circ} \times \Delta, \mathcal{A}) \\
    &K_{**} \ \text{les objets de } \mathcal{A}_{**} \simeq \underline{\mathrm{Hom}}(\Delta^{\circ} \times \Delta^{\circ}, \mathcal{A})
  \end{aligned}
\]
\note{pour $(\mathcal{A}^{*})^{*}$ il écrit $\Delta^{\circ} \times \Delta$, là où l'on attend $\Delta \times \Delta$}

On fera attention que si on compare \add{\uncertain{via}} les foncteurs de symétrie
\struck{$(m, n) \mapsto (n, m)$} \add{$(\Delta^{m}, \Delta^{n}) \mapsto (\Delta^{n}, \Delta^{m})$}
dans $\Delta \times \Delta$ ou dans $\Delta^{\circ} \times \Delta^{\circ}$, ou entre
$\Delta \times \Delta^{\circ}$ et $\Delta^{\circ} \times \Delta$), on n'utilisera pas
la même lettre pour désigner un \struck{complexe} \add{objet} (tel que $K^{*}{}_{*}$
ou $K^{**}$ disons) et son transformé (qui devrait donc s'écrire $K_{*}{}^{*}$, ou
encore une fois $K^{**}$), car cela compliquerait avec la convention plus haut de la
« montée et descente des indices »
$K_{*}{}^{*} \simeq K^{*}{}_{*}{}^{\vee}$ dans le premier cas, et serait
\struck{\ill{}} ambigu : écrire
\[
  K^{[m][n]} = K^{[n][m]}
\]
dans le deuxième ! On utilisera au besoin un symbole « de symétrie » tel que
$\sigma$, pour désigner le symétrique de $K$ ($= K^{*}{}_{*}$, $K^{**}$) par
\[
  {}^{\sigma}K \;(= ({}^{\sigma}K)_{*}{}^{*},\ ({}^{\sigma}K)^{**}) \quad \text{\struck{dans}} \;\ldots
\]
D'ailleurs, \struck{Ces conditions et notations} \add{$\in \mathrm{Ob}\,(\mathcal{A}^{*})^{*}{}_{*}$}
si on a un objet $K^{*}{}_{*}$ disons, alors pour $n \geqslant 0$, on peut considérer
\[
  \text{\struck{\ill{}}} \quad K^{*}{}_{[n]} \overset{\text{df}}{=} (m \longmapsto K^{[m]}{}_{[n]})
\]
comme un objet de $\mathcal{A}^{*}$ ; le foncteur $[n] \longmapsto K^{*}{}_{[n]}$
de $\mathcal{A}^{*}{}_{*}$ n'est autre que ${}^{\sigma}K$.
\note{l'insertion « $\in \mathrm{Ob}\,(\mathcal{A}^{*})^{*}{}_{*}$ » est écrite au-dessus de la ligne biffée, son rattachement est incertain}

On étend ces conventions et notations au cas de symboles $K$ à plus de deux indices,
pour désigner des objets de catégories de

\page{304}
\note{folio de l'auteur (4) ; numéro au crayon des archivistes : 306}
la forme $\mathcal{A}^{*}{}_{*}{}^{**} = [((\mathcal{A}^{*})_{*})^{*}]^{*}$ \struck{\ill{}} p.ex.
--- dans lesquels les objets sont notés \struck{par} des lettres étoilées
$K^{**}{}_{**}{}^{*}$ --- ]
\note{sous le crochet, d'une insertion difficile à placer : « $= \mathrm{Ob}\,\mathcal{S}s^{*}$ »}

Le rôle de $\Delta^{*} \in \mathrm{Ob}\,\widehat{\Delta}^{*}$ comme « objet
cosimplicial universel » est explicité par la

\textbf{Proposition 1.} Soit $\mathcal{A}$ une catégorie où les $\varinjlim$
existent (resp. $\mathcal{B}$ une catégorie où les $\varprojlim$ existent). Alors
\begin{enumerate}
  \item[a)] $\mathcal{A}^{*} \xleftarrow{\ \approx\ } \underline{\mathrm{Hom}}_{\text{comm. à } \varinjlim}(\mathcal{S}s, \mathcal{A})$
  par $F \longmapsto F(\Delta^{*})$
  \item[b)] $\mathcal{B}_{*} \xleftarrow{\ \approx\ } \underline{\mathrm{Hom}}_{\text{comm. à } \varprojlim}(\mathcal{S}s^{\circ}, \mathcal{B})$
  par $F \longmapsto F((\Delta^{*})^{\circ})$
\end{enumerate}
(où $(\Delta^{*})^{\circ}$ désigne l'objet simplicial de $\widehat{\Delta}^{\circ}$
défini par l'objet cosimplicial $\Delta^{*}$ de $\widehat{\Delta}$).

Bien sûr, b) n'est autre que a) appliqué à $\mathcal{A} = \mathcal{B}^{\circ}$ ;
et a) signifie que le foncteur « restriction à $\Delta$ »
\[
  \underline{\mathrm{Hom}}(\widehat{\Delta}, \mathcal{A}) \longrightarrow \underline{\mathrm{Hom}}(\Delta, \mathcal{A})
\]
induit une équivalence sur la s-cat. \struck{$\underline{\mathrm{Hom}}$}
$\underline{\mathrm{Hom}}_{\text{comm. à } \varinjlim}$ du premier membre. (C'est
vrai pour \struck{les} \add{\uncertain{toute}} (petite) catégorie $C$ --- pas
seulement pour $\Delta$, et bien connu.) Le foncteur en sens inverse associe, à
$K_{*} = \varphi : \Delta \to \mathcal{A}$, le foncteur
\[
  X_{*} \longmapsto \varinjlim_{\Delta^{n} \in \Delta/X_{*}} \varphi(\Delta^{n})
  = \varinjlim_{\Delta^{n} \in \Delta/X_{*}} K_{[n]} .
\]
\note{il écrit ici $K_{*}$ et $K_{[n]}$ pour l'objet cosimplicial $\varphi : \Delta \to \mathcal{A}$, là où la notation des pages suivantes demanderait $K^{*}$, $K^{[n]}$ ; l'indice de $\Delta^{n}$ sous les $\varinjlim$ porte une marque illisible}

Dans le cas b), i.e. dans
\[
  \underline{\mathrm{Hom}}(\widehat{\Delta}^{\circ}, \mathcal{B}) \longrightarrow \underline{\mathrm{Hom}}(\Delta^{\circ}, \mathcal{B})
\]
le foncteur en sens inverse associe à
\[
  K^{*} = \psi : \Delta^{\circ} \to \mathcal{B}
\]
le foncteur contravariant
\[
  X^{*} \longmapsto \varprojlim_{\Delta^{n} \in (\Delta/X_{*})^{\circ}} \varphi(\Delta^{n})
  = \varprojlim_{\Delta^{n} \in (\Delta/X_{*})^{\circ}} K^{[n]} .
\]
\note{$\varphi$ et non $\psi$ sous la limite, et $X^{*}$ à gauche pour $X_{*}$, tels quels sur la page}

\textit{Yoga} : nous nous intéressons surtout \struck{aux} (en première étape) aux
invariants covariants attachés à $X_{*}$, à valeurs dans une cat. $\mathcal{A}$,
qui \struck{commutent aux $X_{*}$} \add{commutent aux} $\varinjlim$ par rapport à
$X_{*}$, --- ils sont donc exprimés (si $\mathcal{A}$ est « grande » pour contenir
les $\varinjlim$ voulues) par des éléments de $\mathcal{A}^{*}$, i.e. des objets
\emph{cosimpliciaux} de $\mathcal{A}$. Pour les invariants contravariants en
$X_{*}$, à val. dans une $\mathcal{B}$, \struck{dans} nous abordons en premier lieu
ceux qui \struck{\ill{}} transforment les $\varinjlim$ en $X_{*}$ en $\varprojlim$
dans $\mathcal{B}$ stable par $\varprojlim$ ; ils correspondent \struck{donc} aux
objets de $\mathcal{B}_{*}$, i.e. les objets semi-simpliciaux de $\mathcal{B}$. Si
$\mathcal{B} = (\mathrm{Ens})$, ils correspondent donc à des ens. semi-simpliciaux,
i.e. à des objets $K_{*}$ de $\mathcal{S}s$ \uncertain{lui-même}.
\note{« semi-simpliciaux » : il écrit ½ simpliciaux ; au sens ancien, c'est-à-dire simpliciaux}

\page{305}
\note{numéro au crayon : 307 ; en tête de la colonne droite, « 5 »}
Ce n'est autre que le foncteur représenté par $K_{*}$, et on \add{l'}écrira aussi
\[
  \text{\struck{$K_{*}(X_{*}) = $}} \quad X_{*} \longmapsto K_{*}(X_{*}) = \mathrm{Hom}(X_{*}, K_{*}) .
\]
Si $\mathcal{B} = (\mathcal{A}b)$, $K_{*}$ sera donc un objet de $\mathcal{S}s_{ab}$
\struck{\ill{}}, et les \struck{\ill{}} foncteurs qu'il représente ont
automatiquement \struck{\ill{}} \add{\uncertain{une structure}} \ill{}.
\uncertain{Remarquons} \ill{} les structures algébriques (avec une ou plusieurs
objets de base) qui « peuvent s'exprimer en termes de $\varprojlim$
exclusivement » (groupes, \struck{modules}, anneaux, modules sur
\struck{\ill{}} \add{\uncertain{tels}}, \struck{etc} etc).

\struck{Pour les} \uncertain{Tous} les \struck{foncteurs} invariants ainsi obtenus
à coups de foncteurs représentables --- plus généralement, d'objets simpliciaux ou
cosimpliciaux de catégories $\mathcal{B}$ ou $\mathcal{A}$ --- sont de nature trop
« grosse » et « fruste » pour être intéressants directement --- en particulier, ce
ne sont pas des invariants du type d'homotopie de $X_{*}$. On devra en extraire des
invariants plus subtils qui seront des invariants du type d'homotopie de $X_{*}$,
--- ceux-ci ne seront pas de nature si simple --- i.e. « représentés » par un
objet simplicial ou cosimplicial. Notre propos sera de mettre sur les invariants
« gros » \uncertain{obtenus} \ill{} le \uncertain{dictionnaire}, une machinerie de
structures supplémentaires raisonnables, et dans les catégories de structures de ce
type (complexes de chaînes ou de cochaînes, algèbres cosimpliciales etc)
d'inverser les flèches qui \uncertain{sont} obtenues à partir d'équiv.
d'homotopie dans $\mathcal{S}s$, et de passer à des « catégories de fractions » en
inversant ces flèches. L'invariant ainsi obtenu \struck{\ill{}} \add{plus
« grossier »} bien sûr, mais plus subtil dans sa définition, et moins
« redondant ») exprime donc de façon plus ou moins complète le type d'homotopie,
ou en saisit de façon adéquate tel ou tel aspect.

Dans le cas où $\mathcal{A}$ est une catégorie additive karoubienne (tout
projecteur admet un objet image ou un noyau, ou ce qui revient au même, tt proj. a
un conoyau) la théorie de Dold--Puppe donne des équivalences canoniques
\[
  \mathcal{A}_{*} \;\underset{DP_{*}}{\overset{ND.}{\rightleftarrows}}\; \mathcal{A}. ,
\]
\marginal{$ND$ = « non dégénéré » \quad $DP$ = « Dold--Puppe »}
\note{les deux flèches portent chacune le signe $\approx$}
où $\mathcal{A}.$ désigne la catégorie des \emph{complexes de chaînes} dans
$\mathcal{A}$. On désigne par une \add{telle} \struck{=} lettre \add{$K$} deux
objets qui se correspondent, mais en mettant un indice bas $*$ ou bas $\cdot$,
suivant qu'il

\page{306}
s'agit de désigner l'avatar « objet simplicial » ou « complexe de chaînes » :
\[
  K. = ND.(K_{*}), \qquad K_{*} = DP_{*}(K.)
\]
\[
  K_{*} = (n \mapsto K_{[n]}), \qquad K. = (K_{n}, d_{n})_{n \geqslant 0}
\]
On est ainsi amené à noter par une notation $K.$ tel complexe de chaînes, par $K_{n}$
ses composantes (ce qui exclut l'utilisation de la notation $K_{n}$ pour la
$n$-ième composante \add{de} $K_{*}$, et nous oblige à adopter une notation telle
que $K_{[n]}$ !).

De même, la théorie de DP \uncertain{appliquée} à $\mathcal{A}^{\circ}$ nous donne
des équivalences
\[
  \mathcal{A}^{*} \;\underset{DP^{*}}{\overset{ND^{\cdot}}{\rightleftarrows}}\; \mathcal{A}^{\cdot}
\]
et des notations
\[
  K^{\cdot} = ND^{\cdot}(K^{*}), \qquad K^{*} = DP^{*}(K^{\cdot})
\]
\[
  K^{*} = (n \mapsto K^{[n]}), \qquad K^{\cdot} = (K^{n}, d^{n})_{n \geqslant 0} .
\]
Les foncteurs $ND^{\cdot}$ et $DP^{*}$ commutent aux foncteurs additifs. En
particulier, dans le cas d'une équivalence
\[
  \mathcal{A} \xleftrightarrow[\ v\ ]{\ \approx\ } \mathcal{B}^{\circ}
\]
dite \uncertain{dualité}, étendant aux complexes de chaînes ou de cochaînes la
convention de la montée ou de la descente des indices, on trouve que les foncteurs
$ND$ et $DP$ y commutent. Donc dans ce cas, la même lettre $K$ désignera

\page{307}
quatre avatars possibles
\[
  K^{*} \quad K_{*} \quad K^{\cdot} \quad K.
\]
liés par les relations ci-dessus
\[
  \begin{cases}
    K^{*} = DP^{*}(K^{\cdot}) & K^{\cdot} = ND^{\cdot}(K^{*}) \\
    K_{*} = DP_{*}(K.) & K. = ND.(K_{*})
  \end{cases}
\]
plus les relations
\[
  \begin{cases}
    K^{*} = (K_{*})^{\vee}, & K_{*} = (K^{*})^{\vee} \\
    K^{\cdot} = (K.)^{\vee}, & K. = (K^{\cdot})^{\vee} .
  \end{cases}
\]
\struck{(\ill{})} Ces \struck{$K$} notations \uncertain{est} vaguement (i.e. à
équivalence de catégories) s'étendent aux objets \struck{\ill{}} \add{dépendant de}
plusieurs indices, p.ex. \struck{$K^{*}$} une même lettre $K$ (désignant
\add{\uncertain{p.ex. initialement}} un objet $K^{*}{}_{*}$ de
$(\mathcal{A}_{*})^{*}$, \add{disons}) désigne également \struck{\ill{}} quatre
avatars distincts
\[
  \begin{cases}
    K^{*}{}_{*} \in (\mathcal{A}_{*})^{*}, & (K^{\cdot}{}_{*}) \in (\mathcal{A}_{*})^{\cdot} \\
    K^{\cdot}{}_{\cdot} \in (\mathcal{A}.)^{\cdot}, & (K^{*}{}_{\cdot}) \in (\mathcal{A}.)^{*}
  \end{cases}
\]
\note{dans $(\mathcal{A}_{*})^{\cdot}$ l'exposant est surchargé (point sur une étoile)}
avec des notations
\[
  K^{[n]}{}_{[m]}, \quad K^{n}{}_{[m]}, \qquad K^{n}{}_{m}, \quad K^{[n]}{}_{m}
\]
pour les composantes. De même, on aura, si on part p.ex. de
$K^{**} \in (\mathcal{A}^{*})^{*}$,
\[
  \begin{cases}
    K^{**} \in (\mathcal{A}^{*})^{*}, & K^{\cdot *} \in (\mathcal{A}^{\cdot})^{*} \\
    K^{\cdot\cdot} \in (\mathcal{A}^{\cdot})^{\cdot}, & K^{*\cdot} \in (\mathcal{A}^{*})^{\cdot}
  \end{cases}
\]
\[
  \begin{cases}
    K^{[n][m]}, & K^{n[m]} \\
    K^{nm}, & K^{[n]m}
  \end{cases}
\]
On notera quand objets de $(\mathcal{A}_{*})^{\cdot}$ ou un complexe de
cochaînes \add{$n \mapsto (K_{*}{}^{n}, d^{n})$} dans la catégorie $\mathcal{A}_{*}$
des objets semi-simpliciaux de $\mathcal{A}$,
\note{« quand » tel quel ; l'insertion $n \mapsto (K_{*}{}^{n}, d^{n})$ est en interligne, sa fin lue avec doute}

\page{308}
\ill{} au dessus (moyennant \uncertain{l'}équivalence de catégories canonique
\[
  \sigma : (\mathcal{A}_{*})^{\cdot} \cong (\mathcal{A}^{\cdot})_{*} \;)
\]
un objet semi-simplicial dans la catégorie $\mathcal{A}^{\cdot}$ des complexes de
cochaînes, noté
\[
  n \longmapsto K_{[n]}{}^{\cdot} .
\]
On interprète de même les objets des autres catégories mixtes
\struck{\ill{}} :
\[
  (*) \quad
  \begin{cases}
    \mathcal{A}_{*\cdot}, \ \mathcal{A}_{\cdot *}, \ \mathcal{A}^{**}{}_{\cdot}, \ \text{\struck{\ill{}}} \\
    \mathcal{A}^{*}{}_{*}, \ \mathcal{A}^{\cdot}{}_{*}, \ \mathcal{A}^{\cdot}{}_{*}, \ \text{\struck{\ill{}}}
  \end{cases}
\]
\note{les indices des catégories mixtes de la liste $(*)$ sont surchargés et lus avec doute ; la fin de chaque ligne est biffée en hachures}
De façon analogue, les objets $K^{\cdot\cdot}$ de $(\mathcal{A}^{\cdot})^{\cdot}$
sont les bicomplexes de cochaînes, les objets $K..$ de
$(\mathcal{A}.). = \mathcal{A}..$ sont les bicomplexes de chaînes, enfin les objets
$K^{\cdot}.$ de $(\mathcal{A}.)^{\cdot}$ ou $K.^{\cdot}$ de $(\mathcal{A}^{\cdot}).$
sont les bicomplexes mixtes \add{chaînes--cochaînes} (de chaînes par un indice, de
cochaînes par un autre). Dans le cas de $K^{\cdot\cdot}$ ou $K..$ le complexe simple
associé est un complexe de cochaînes resp. de chaînes, tandis que dans le cas mixte
$K^{\cdot}.$ ou $K.^{\cdot}$, on trouve des complexes simples avec des objets en
degrés positifs et négatifs à priori (il y aurait lieu de fixer une convention de
degré, si pour le degré total l'\uncertain{opérateur} différentiel est de degré
$+1$ ou $-1$, suivant

\page{309}
qu'on donne la préférence \struck{\ill{}} cochaînes ou aux chaînes. (Je n'ai pas eu
encore besoin \emph{vraiment} d'une telle convention de notations…)

Dans le cas d'une dualité
\[
  v : \mathcal{A} \xrightarrow{\ \approx\ } \mathcal{B}^{\circ}
\]
on étend les conventions de montée et descente des indices aux objets à plusieurs
indices, ainsi
\[
  (K_{*}{}^{\cdot})^{\vee} = K^{*}{}_{\cdot} .
\]
\note{la flèche de $v$ est doublée d'un trait vers la gauche ; à droite de la formule, une marque biffée}
Dans la suite d'objets mixtes $(*)$ de la page précédente, on a écrit l'un au-dessous
de l'autre des objets qui se correspondent par dualité.

\textit{Remarque.} Il est essentiel pour notre système de conventions de définir
(disons) des bicomplexes de cochaînes (objets de $\mathcal{A}^{\cdot\cdot}$) avec
deux \uncertain{op. de diff.} \ill{}, qui \emph{commutent} --- on n'utilisera donc
pas les conventions de Cartan--Eilenberg. Il faudra donc un signe dans la diff.
totale.

Du point de vue de la page 1 et du yoga \struck{\ill{}} \ill{} « construction
d'invariants associés à des $X_{*} \in \mathrm{Ob}\,\mathcal{S}s$ », on peut dire que
la théorie de Dold--Puppe

\page{310}
\note{folio de l'auteur 10 ; numéro au crayon : 312}
fournit
\struck{\ill{}}

\textbf{Prop. 2.} Avec les notations \add{et hyp.} de prop. 1, avec $\mathcal{A}$
et $\mathcal{B}$ additives \struck{\ill{}}
\begin{enumerate}
  \item[a)] $\underline{\mathrm{Hom}}_{\text{comm. aux } \varinjlim}(\mathcal{S}s, \mathcal{A}) \xrightarrow{\ \approx\ } \mathcal{A}^{\cdot}$
  \item[b)] $\underline{\mathrm{Hom}}_{\text{comm. aux } \varinjlim}((\mathcal{S}s)^{\circ}, \mathcal{B}) \xrightarrow{\ \approx\ } \mathcal{B}.$
\end{enumerate}
\note{sous les deux $\underline{\mathrm{Hom}}$, l'indice abrégé lu comme « comm. aux $\varinjlim$ », avec doute ; en b) on attendrait $\varprojlim$, la flèche est ambiguë}

Nous allons \struck{préciser} \add{expliciter} cette correspondance du point de vue
de l'algèbre universelle. Prenons d'abord le cas de b), avec $\mathcal{B} = (\mathcal{A}b)$
pour simplifier. \struck{\ill{}} \add{Dans \uncertain{cette} \uncertain{optique}}
précisons la relation entre
\[
  F : (\mathcal{S}s)^{\circ} \xrightarrow[\text{aux } \varinjlim]{\text{commutant}} \mathcal{A}b
  \quad \text{et} \quad K. = (K_{n}, d_{n}) \in \mathcal{A}b.
\]
qui se déterminent mutuellement. On désigne comme de juste par $K_{*}$ l'objet de
$\mathcal{A}b_{*}$ défini par $K.$ On a donc
\[
  F(X_{*}) \simeq \mathrm{Hom}_{\widehat{\Delta}}(X_{*}, K_{*})
  \simeq \mathrm{Hom}_{\widehat{\Delta}_{ab}}(\mathbb{Z}(X_{*}), K_{*})
\]
où $\mathbb{Z}(X_{*})$ désigne le « faisceau abélien libre engendré par $X_{*}$ »,
ou si on veut
\[
  \mathbb{Z}(X_{*}) = (n \longmapsto \mathbb{Z}(X_{n}) = \mathbb{Z}^{(X_{n})})
\]
Ce groupe \add{ab.} semi-simplicial sera aussi noté $\mathbb{C}_{*}(X_{*})$, et
$\mathbb{C}.(X_{*})$ le complexe de chaînes associé \add{par DP} :
\struck{on aura donc}
\[
  \begin{cases}
    \mathbb{C}_{*}(X_{*}) = \mathbb{Z}(X_{*}) = (n \longmapsto \mathbb{Z}(X_{n}) = \mathbb{Z}^{(X_{n})}) \\
    \mathbb{C}.(X_{*}) = ND.(\mathbb{C}_{*}(X_{*}))
  \end{cases}
\]
\note{« par DP » est en interligne, alors que la formule écrit $ND.$}
Donc par Dold--Puppe on aura
\[
  F(X_{*}) \simeq \mathrm{Hom}_{\widehat{\Delta}_{ab} = \mathcal{A}b_{*}}(\mathbb{C}_{*}(X_{*}), K_{*})
  \underset{\text{D.P.}}{\simeq} \mathrm{Hom}_{\mathcal{A}b.}(\mathbb{C}.(X_{*}), K.)
\]
Donc le foncteur $F_{K.} = F : \mathcal{S}s^{\circ} \longrightarrow \mathcal{A}b$
\[
  X_{*} \longmapsto F_{K.}(X_{*}) = F(X_{*})
\]
associé à $K. \in \mathcal{A}b.$ est donné par
\[
  \boxed{F_{K.}(X_{*}) = \mathrm{Hom}.(\mathbb{C}.(X_{*}), K.)}
\]
On en tire deux conclusions :

1°) Les valeurs sur $X_{*}$ des foncteurs envisagés $F$ ($\mathcal{S}s^{\circ} \to \mathcal{A}b$,
commutant aux $\varinjlim$) ne dépendent que de \add{(« abélianisé de $X_{*}$ »),}
$\mathbb{C}_{*}(X_{*}) \in \mathcal{A}b_{*}$, ou ce qui revient au même, de
$\mathbb{C}.(X_{*}) \in \mathcal{A}b.$ ; \struck{\ill{}} ce complexe de chaînes
$\mathbb{C}.(X_{*})$, peut être considéré comme « l'invariant additif universel » de
$X_{*}$. Quand on « oublie $X_{*}$ au profit de $\mathbb{C}.(X_{*}) = C. \in \mathcal{A}b.$ »
les foncteurs envisagés (en tant que foncteurs en $C.$) sont exactement les
foncteurs représentables sur $\mathcal{A}b.$.

[\emph{NB}. Quand le foncteur $F$ a une structure plus riche, p.ex. une
multiplication $F \otimes F \to F$ i.e. est à valeurs dans les $\mathbb{Z}$-algèbres,
\struck{et \ill{}} cette structure sur $F(X_{*})$ n'est pas reconstruite par

\page{311}
\note{folio de l'auteur 11 ; numéro au crayon : 313}
la seule connaissance de $C. \in \mathcal{A}b.$, au lieu de $X_{*}$ ; il faut, pour
\struck{\ill{}} \add{\uncertain{les}} décrire en \emph{termes de} $C.$ \add{certaines
des} structures \add{ainsi} qui peuvent \uncertain{exister} sur $F(X_{*})$, en termes
de celle de $F$, de tenir compte de structures supplémentaires éventuelles de $C.$,
qui peuvent s'expliciter en se rappelant que $C.$ provient d'un $X_{*}$.
\struck{\ill{}} \add{\uncertain{Cela} \uncertain{peut} être considéré comme}
\struck{\ill{}} d'étude systématique de ces structures \struck{\ill{}}
\struck{\ill{}} (le contenu essentiel de nos réflexions).

Mais bien sûr, \struck{si les structures \ill{}} \add{dans le} \uncertain{cas} où
on a sur $F$ une structure de $k$-module, i.e. $F : \mathcal{S}s^{\circ} \to \mathcal{A}b_{k}$
(où $\mathcal{A}b_{k}$ est la catégorie des $k$-modules), \struck{il suffit}
\add{\uncertain{il est possible}} de se donner la structure \uncertain{simpl.} sur
$C.$, ou \struck{\ill{}} qui en fait \struck{de complexe de \ill{} \ill{} \ill{}
(i.e. se \ill{} $k \to \mathrm{End}_{\mathcal{A}b}(C.)$) \ill{} récupérer la
structure complète de} par les $K.$ qui décrivent $F$ maintenant dans
$\mathcal{A}b_{k}.$ i.e. par
\[
  k \longrightarrow \mathrm{End}_{\mathcal{A}b.}(K.)
\]
la \uncertain{qui} \uncertain{est} aussi sur $\mathrm{Hom}_{\mathcal{A}b.}(C., K.)$,
dont on récupère ainsi la structure de $k$-module. En fait il suffit, au lieu de
$C.$, de connaître \uncertain{seulement} $C_{k\cdot} = C. \otimes_{\mathbb{Z}} k$
(noté $\mathbb{C}_{k*}(X_{*})$), car on aura
\note{un passage encadré et barré de traits obliques, au milieu de la colonne gauche, est biffé ; on n'en donne que ce qui se lit}

\[
  \text{\struck{\ill{}}} \quad
  \begin{cases}
    F(X_{*}) \simeq \mathrm{Hom}_{(\mathcal{A}b_{k}).}(C_{k\cdot}, K.) \\
    C_{k\cdot} = \mathbb{C}_{k*}(X_{*}) = \mathbb{C}.(X_{*}) \otimes_{\mathbb{Z}} k
  \end{cases}
\]
Plus généralement, revenant au cas d'une $\mathcal{B}$ cat. additive générale
stable par $\varinjlim$ et $F : (\mathcal{S}s)^{\circ} \to \mathcal{B}$ et
$K. \in \mathcal{B}.$ se correspondant par la prop.~2, on peut encore écrire
\[
  F(X_{*}) \simeq \mathrm{Hom}^{\cdot}.(\mathbb{C}.(X_{*}), K.)
\]
en définissant la deuxième expression comme \emph{objet de $\mathcal{B}$}, en
définissant d'abord le bicomplexe de $\mathcal{B}$
\[
  \mathcal{H}om^{\cdot}.(C., K.) = (\mathrm{Hom}_{\mathcal{B}}(C_{i}, K_{j}))_{i,j}
\]
(donc $\mathcal{H}om^{i}{}_{j}(C., K.) = \mathrm{Hom}_{\mathcal{B}}(C_{i}, K_{j})$)
--- où l'objet $\mathrm{Hom}(A, M)$, pour $A \in \mathcal{A}b$ et $M \in \mathcal{B}$,
se définissant par
\[
  \mathrm{Hom}_{\mathcal{B}}(N, \mathrm{Hom}(A, M)) \simeq \mathrm{Hom}_{\mathcal{A}b}(A, \mathrm{Hom}_{\mathcal{B}}(N, M))
\]
(le foncteur \struck{\ill{}} \add{du 2\supplied{e}} \struck{\ill{}} étant bien
représentable en $N$, pour $A, M$ fixés, grâce à l'hyp. faite sur $\mathcal{B}$
d'être stable par $\varprojlim$).
On définit alors
\[
  \mathrm{Hom}.(\mathbb{C}.\text{\struck{\ill{}}}, K.) \overset{\text{df}}{=} Z^{0}(\mathcal{H}om^{\cdot}.(C., K.))
\]
\add{cycles de degré $0$ pour degré total (\uncertain{associé} au bicomplexe)}
Si enfin $\mathcal{B}$ est une cat. $k$-linéaire ($k$ commutatif) i.e. munie d'un
homom. d'anneaux $k \to \mathrm{End}(\mathrm{id}_{\mathcal{B}})$,
\note{« stable par $\varinjlim$ » plus haut, « stable par $\varprojlim$ » ici, tels quels}

\page{312}
\note{folio de l'auteur 12 ; numéro au crayon : 314}
alors la connaissance de $C_{k\cdot} = \mathbb{C}_{k\cdot}(X_{*}) = C. \otimes_{\mathbb{Z}} k$
est encore suffisante pour la construction de $F(X_{*})$, car on aura maintenant
\[
  F(X_{*}) \simeq \mathrm{Hom}_{k.}(C._{k}, K.)
  \simeq Z^{0}(\underbrace{\mathcal{H}om_{k}{}^{\cdot}.(C._{k}, K.)}_{\mathcal{H}om_{k}{}^{i}{}_{j}(C._{k}, K.) = \mathcal{H}om_{k}(C_{i,k}, K_{j})})
\]
($\mathcal{H}om_{k}(A, M)$, pour $A \in \mathcal{A}b_{k}$ et $M \in \mathcal{B}$, étant défini par
\[
  \mathrm{Hom}_{\mathcal{B}}(N, \mathcal{H}om_{k}(A, M)) \simeq \mathrm{Hom}_{k}(A, \mathrm{Hom}_{\mathcal{B}}(N, M))).
\]

2°) Si on regarde toutes les façons possibles \struck{d'\ill{}} \add{d'envoyer}
\[
  (\mathcal{S}s)^{\circ} \xrightarrow{\ F\ } \mathcal{B}
\]
dans $\mathcal{B}$ stable par $\varprojlim$, avec $F$ y commutant, il y en a
\struck{\ill{}} une « universelle », dont toutes les autres sont « spécialisation »,
obtenue en prenant $\mathcal{B}_{0} = (\mathcal{A}b.)^{\circ}$, et
$F_{0} : (\mathcal{S}s)^{\circ} \to \mathcal{B}_{0} = (\mathcal{A}b.)^{\circ}$ défini par
\[
  F_{0}(X_{*}) = \mathbb{C}.(X_{*})
\]
(qui dépend de $X_{*}$ de façon contravariante quand $\mathbb{C}.(X_{*})$ s'envisage
comme objet de $(\mathcal{A}b.)^{\circ}$). De façon précise ;

\marginal{\uncertain{Nouvelle} forme de l'énoncé du théorème de Dold--Puppe}

\textbf{Prop. 3.} Considérons
\[
  \mathcal{S}s^{\circ} \xrightarrow{\ F_{0}\ } \mathcal{B}_{0} = (\mathcal{A}b.)^{\circ} \simeq (\mathcal{A}b^{\circ})^{\cdot}
\]
défini par $F_{0}(X_{*}) = \mathbb{C}.(X_{*})$. Alors $\mathcal{B}_{0}$ est une
catégorie \add{add.} avec $\varprojlim$, et pour toute autre telle catégorie
$\mathcal{B}$, le foncteur
\[
  \underline{\mathrm{Hom}}_{\text{comm. aux } \varprojlim}(\text{\struck{\ill{}}}\,\mathcal{B}_{0}, \mathcal{B})
  \longrightarrow \underline{\mathrm{Hom}}_{\text{comm. } \varprojlim}(\mathcal{S}s^{\circ}, \mathcal{B})
  \quad (\simeq \mathcal{B}. \text{ par prop. 2})
\]
donné par $\varphi \longmapsto \varphi \circ F_{0}$ est une équivalence de catégories
\struck{\ill{} \ill{} prop. 2, cela signifie \ill{} \ill{} grâce au théorème de
Dold--Puppe, \ill{} \ill{} par $F_{0}$ \ill{} \ill{} grâce \ill{} : l'équivalence
\ill{} \ill{} \ill{} indique la condition \ill{} aux $\varprojlim$)}
\note{un passage de six lignes environ, encadré et barré de hachures, est biffé ; on n'en donne que les mots lisibles}
$\underline{\mathrm{Hom}}_{!}(\mathcal{B}_{0}, \mathcal{B})$ remplaçons dans la déf. de
$F_{0}$, $\mathcal{B}_{0}$, $(\mathcal{A}b.)^{\circ}$ par
$(\mathcal{A}b_{*})^{\circ} = \widehat{\Delta}_{ab}{}^{\circ}$, $\mathbb{C}.(X_{*})$ par
$\mathbb{C}_{*}(X_{*})$ \add{$= \mathbb{Z}(X_{*})$} et l'équivalence s'écrit
\[
  \underline{\mathrm{Hom}}_{!}(\widehat{\Delta}_{ab}{}^{\circ}, \mathcal{B})
  \xrightarrow{\ \approx\ } \underline{\mathrm{Hom}}_{!}(\widehat{\Delta}^{\circ}, \mathcal{B})
\]
où l'exposant $!$ signifie qu'on se borne aux foncteurs qui commutent aux $\varprojlim$.
\note{sous les deux $\underline{\mathrm{Hom}}$ de la formule, un indice surchargé, illisible}
Or sous cette forme, l'assertion est valable quelle que soit la cat. (non pointée)
$\Delta$, [\add{on a un foncteur « abélianisation »} \struck{\ill{} \ill{} se faisant
\ill{}}
\[
  X \longmapsto \mathbb{Z}(X.) \qquad \widehat{\Delta} \longrightarrow \widehat{\Delta}_{ab}
\]
\struck{commutant aux $\varinjlim$ \ill{}} \struck{$\widehat{\Delta}^{\circ} \to \mathcal{B} = (\widehat{\Delta}_{ab})^{\circ}$}

\page{313}
\note{folio de l'auteur 13 ; numéro au crayon : 315}
donc
\[
  \widehat{\Delta}^{\circ} \xrightarrow{\ F_{0}\ } \mathcal{B}_{0} = (\widehat{\Delta}_{ab})^{\circ}
  \cong \underline{\mathrm{Hom}}(\Delta^{\circ}, \mathcal{A}b)^{\circ}
  \approx \underline{\mathrm{Hom}}(\Delta, \mathcal{A}b^{\circ})
\]
\[
  [\approx (\mathcal{A}b^{\circ})^{*} \approx (\mathcal{A}b_{*})^{\circ} \text{ dans le cas qui nous intéresse}]
\]
commutant aux $\varinjlim$ d'où le foncteur
\[
  \underline{\mathrm{Hom}}_{!}(\text{\struck{\ill{}}}\,\mathcal{B}_{0}, \mathcal{B}) \longrightarrow \underline{\mathrm{Hom}}_{!}(\widehat{\Delta}^{\circ}, \mathcal{B})
\]
\struck{et on définit un foncteur} dont on veut montrer que c'est une équiv. Posant
$\mathcal{B} = \mathcal{A}^{\circ}$, $\mathcal{B}_{0} = (\mathcal{A}_{0})^{\circ}$, et
$F_{0} = (G_{0})^{\circ}$, donc
\[
  G_{0} : \widehat{\Delta} \to \widehat{\Delta}_{ab}, \qquad G_{0}(X) = \mathbb{Z}(X),
\]
on doit prouver que
\[
  \underline{\mathrm{Hom}}_{!}(\widehat{\Delta}_{ab}, \mathcal{A}) \longrightarrow \underline{\mathrm{Hom}}_{!}(\widehat{\Delta}, \mathcal{A}),
  \qquad \varphi \longmapsto \varphi \circ G_{0}
\]
est une équivalence, où $!$ désigne les foncteurs qui commutent aux $\varinjlim$. Or le
deuxième membre est $\approx \underline{\mathrm{Hom}}(\Delta, \mathcal{A})$ (par
restriction à $\Delta \hookrightarrow \widehat{\Delta}$), \struck{donc} le premier
s'écrit ? \struck{on} Prenons $\varphi \in \underline{\mathrm{Hom}}_{!}(\widehat{\Delta}_{ab}, \mathcal{A})$
et notons que pour tout $M \in \mathcal{A}$, \struck{l'\ill{}} le foncteur
$X \mapsto \mathrm{Hom}(\varphi X, M)$ transforme $\varinjlim$ en $\varprojlim$, donc
est représentable
\[
  \mathrm{Hom}(\varphi X, M) \simeq \mathrm{Hom}(X, \psi M)
\]
i.e. l'hyp. de commutation de $\varphi$ aux $\varinjlim$ s'exprime par le fait qu'il
admet un adjoint $\psi$ --- la condition \uncertain{ici} \uncertain{mise} sur
\struck{\ill{}} $\psi : \mathcal{A} \to \widehat{\Delta}_{ab}$ pour qu'il provienne
d'un $\varphi$ étant que pour tout $X$ fixé dans $\widehat{\Delta}_{ab}$,
$M \mapsto \mathrm{Hom}(X, \psi M)$ soit repr. Écrivant
$\widehat{\Delta}_{ab} \simeq \underline{\mathrm{Hom}}(\Delta^{\circ}, (\mathcal{A}b))$
en termes d'interprétation de $\psi$ \uncertain{comme} cat.
$\underline{\mathrm{Hom}}_{!}(\widehat{\Delta}_{ab}, \mathcal{A})$ des $\varphi$ comme
\uncertain{sous}-cat. pleine de $\underline{\mathrm{Hom}}(\mathcal{A} \times \Delta^{\circ}, \mathcal{A}b)$
\add{$\simeq \underline{\mathrm{Hom}}(\Delta^{\circ}, \underline{\mathrm{Hom}}(\mathcal{A}, \mathcal{A}b))$
\uncertain{sous-cat. des}} foncteurs $\lambda$ satisfaisant une condition de repr. \add{ci-dessus} en
$\mathcal{A}$. Or si $x \in \Delta$, le foncteur \struck{\ill{}}
\[
  \lambda(x) : \mathcal{A} \to \mathcal{A}b \ (\in \mathcal{A}b \ldots)
\]
\note{après $\lambda$, deux arguments surchargés et biffés ; l'insertion « $\lambda(x) : \mathcal{A} \to$ » est lue avec doute}
n'est autre que
\[
  M \longmapsto \mathrm{Hom}(\mathbb{Z}(x), \psi(M))
\]
(car
\[
  \lambda(x)(M) \overset{\text{df}}{=} \psi(M)(x) \simeq \mathrm{Hom}_{\widehat{\Delta}}(x, \psi(M))
  = \mathrm{Hom}_{\widehat{\Delta}_{ab}}(\mathbb{Z}(x), \psi(M))
\]
donc il doit être représentable, \struck{\ill{}} \uncertain{grâce} \uncertain{à}
$X = \mathbb{Z}(x)$) donc \uncertain{par} $\varphi(x) \in \mathcal{A}$, donc on est
dans $\underline{\mathrm{Hom}}(\Delta^{\circ}, \text{\struck{\ill{}}}\,\mathcal{A}^{\circ}) \cong \underline{\mathrm{Hom}}(\Delta, \mathcal{A})$,
on a pratiquement terminé…]

Quel est \add{donc} le foncteur commutant aux $\varinjlim$ \uncertain{le plus}
universel de $\widehat{\Delta}$ dans une cat. additive $\mathcal{A}$ avec $\varinjlim$ ?
C'est
\[
  X \longmapsto \mathbb{Z}(X) = \mathbb{Z}^{(X)} \quad \text{de } \widehat{\Delta} \text{ dans } \widehat{\Delta}_{ab}
\]
(ou le foncteur « \emph{abélianisation} ») \uncertain{sans} doute se \uncertain{vérifie}
pour un topos quelconque. \add{(i.e. $\Delta$ = cat. des $\Delta^{n}$)}
Mais ici \add{cela signifie}

\emph{Variante} équiv. \struck{\ill{}} \add{(a)} (Dold--Puppe) : l'objet cosimplicial
\struck{\ill{}} $n \longmapsto \mathbb{C}_{*}(\Delta^{n}{}_{*}) = \mathbb{Z}(\Delta^{n}{}_{*}) = \mathbb{Z}^{\Delta^{n}{}_{*}}$
de $\widehat{\Delta}_{ab} \simeq \mathcal{A}b_{*}$ est \emph{universel} pour les objets
cosimpliciaux dans des cat. additives avec $\varinjlim$ quelc. (pour des foncteurs entre
telles cat. qui commutent aux $\varinjlim$). Donc l'\struck{objet} \add{complexe de
cochaînes correspondant}

\page{314}
\note{folio de l'auteur 14 ; numéro au crayon : 316}
dans $\widehat{\Delta}_{ab} \simeq \mathcal{A}b_{*}$ est \emph{universel} pour les
complexes de cochaînes dans des cat. additives $\mathcal{A}$ avec $\varinjlim$. \ill{}
(\emph{NB} \uncertain{On} \uncertain{peut} aussi interpréter $\mathcal{A}b_{*}$ comme
$\mathcal{A}b.$ dans cette description de la cat. contenant l'objet universel…).

b) En dualisant et interprétant $n \longmapsto \mathbb{C}_{*}(\Delta^{n}{}_{*})$ comme
objet \emph{simplicial} dans $(\mathcal{A}b_{*})^{\circ}$ \add{$= (\mathcal{A}b^{\circ})^{*}$},
\struck{\ill{}} regardée comme cat. additive avec $\varprojlim$, il est universel pour
les cat. $\mathcal{B}$ avec $\varprojlim$ et foncteurs entre icelles y commutant. De
même pour le complexe de chaînes associé.

Du point de vue de l'algèbre universelle, \struck{on \ill{} que les \ill{}}
\add{on peut exprimer le contenu de ces} réflexions de la façon suivante :
\begin{enumerate}
  \item[1.] Les inv. contravariants « commutant aux $\varprojlim$ » en $X_{*} \in \mathcal{S}s$,
  à valeurs dans une cat. \emph{additive} $\mathcal{A}$ avec $\varprojlim$, s'expriment
  en termes d'objets simpliciaux ou en termes de complexes de chaînes de $\mathcal{A}$.
  \item[2.] Les constructions qui se font sur de tels invariants \add{(\uncertain{pourvu} qu'elles
  \uncertain{soient}} \struck{\ill{}} \uncertain{compatibles} aux foncteurs
  $\mathcal{A} \to \mathcal{A}'$ commutant aux $\varprojlim$) s'expriment : à l'aide des
  objets de $(\mathcal{A}b_{*})^{\circ}$ (où se trouve l'objet simplicial --- pour cat.
  add. à $\varprojlim$ --- \emph{universel}), ou de $(\mathcal{A}b.)$
\end{enumerate}
et $1'$), $2'$) les énoncés duals (formels à partir de 1), 2)).

Si $F$, $K_{*}$, $K.$ se correspondent dans 1, on a
\[
  \begin{cases}
    F(X_{*}) \simeq \mathrm{Hom}_{*}(\mathbb{C}_{*}(X_{*}), K_{*}) \simeq \mathrm{Hom}.(\mathbb{C}.(X_{*}), K.) \\
    K. = ND.(K_{*}), \quad K_{*} = DP_{*}(K.)
  \end{cases}
\]
Si de même $\Omega$ est une « opération sur des invariants », et si $L_{*}$ \add{$\in \mathcal{A}b_{*}$}
et $L. \in \mathcal{A}b.$ y sont associés, on aura
\[
  \begin{cases}
    \Omega(F) \simeq \mathrm{Hom}_{*}(L_{*}, K_{*}) \simeq \mathrm{Hom}.(L., K.)
    \qquad (= \Omega_{*}(K_{*}),\ = \Omega.(K.)) \\
    L. = ND.(L_{*}), \quad L_{*} = DP_{*}(L.)
  \end{cases}
\]
\marginal{$\Omega$ a deux expressions « opération sur des $K^{*}$ » ($\Omega^{*}$) et « op. sur des $K^{\cdot}$ » ($\Omega^{\cdot}$) --- opérations sur des objets de $\mathcal{B}$ \struck{\ill{}} : notion d'un complexe (semi-simpl. ou de chaînes) dans $\mathcal{B}$}
(La signification des symboles $\mathrm{Hom}_{*}$ \add{des $\mathcal{A}b_{*}$} (resp.
$\mathrm{Hom}.$ \add{des $\mathcal{A}b.$}) pour \struck{\ill{}} un $C_{*}$ ou $L_{*}$
(resp. $C.$ ou $L.$) a été expliquée plus haut).

Dualement, \add{pour} les \struck{\ill{}} \add{invariants} \add{d'objets $X_{*} \in \mathcal{S}s$}
covariants commutant aux $\varinjlim$, à valeurs dans $\mathcal{A}$, si $G$, $K^{*}$,
$K^{\cdot}$ se correspondent, on aura
\[
  \begin{cases}
    G(X_{*}) \simeq \text{\struck{\ill{}}}\;\mathbb{C}_{*}(X_{*}) \overset{*}{\otimes} K^{*} \simeq \text{\struck{\ill{}}}\;\mathbb{C}.(X_{*}) \dot{\otimes} K^{\cdot} \\
    K^{\cdot} = ND^{\cdot}(K^{*}), \quad K^{*} = ND^{*}(K^{\cdot})
  \end{cases}
\]
\marginal{: préciser la signification de ces formules}
\note{$K^{*} = ND^{*}(K^{\cdot})$ tel quel, là où l'on attend $DP^{*}$}
et si l'opération $\Omega$, $L_{*}$ et $L.$ se correspondent, ($L_{*} \in \mathcal{A}b_{*}$,
$L. \in \mathcal{A}b.$) on aura
\[
  \begin{cases}
    \Omega(G) \simeq L_{*} \overset{*}{\otimes} K^{*} \simeq L. \dot{\otimes} K^{\cdot} \\
    L. \simeq ND.(L_{*}), \quad L_{*} \simeq DP_{*}(L.)
  \end{cases}
\]

\page{315}
\note{folio de l'auteur 15 ; numéro au crayon : 317}
\emph{Remarques.} a) On trouve en utilisant les \uncertain{deux} objets ($L_{*}$ ou
$L.$) pour exprimer des \emph{opérations} sur invariants additifs contravariants, ou sur
invariants \add{« additifs »} covariants (du type envisagé) --- c'était clair à priori
à cause du passage de $\mathcal{A}$ à $\mathcal{B}$ par \add{passage à la} catégorie
opposée --- mais l'op. définie par $L_{*}$ \struck{\ill{}} et $L.$ de façon
contravariante dans le premier cas, covariante dans le deuxième.

b) On a en principe résolu, et de façon quasi-tautologique, la question initiale
des \struck{\ill{} \ill{}} \add{déterminations des constructions} possibles sur
invariants contravariants ou covariants « additifs » du type envisagé : on trouve
\struck{\ill{}} \emph{\uncertain{grosses}} catégories, $\mathcal{A}b_{*}$ ou
$\mathcal{A}b.$ (ou \struck{\ill{}} passage qui de celles-ci à l'opposée). Cela
tient au fait que l'on a admis des catégories de valeurs \add{$\mathcal{A}$}
\add{$\mathcal{B}$ ou $\mathcal{A}$,} assez générales \struck{\ill{}} où on
peut effectuer \struck{\ill{}} les $\varprojlim$, ou bien les $\varinjlim$ --- ce
qui a pour effet que un \uncertain{tas} d'opérations peuvent s'effectuer sur les
\struck{\ill{}} \struck{invariants} invariants, et op. sur les invariants.

On trouvera par \uncertain{lors} qu'en est plus restrictif sur les opérations qu'on
se permet, \struck{\ill{}} \add{\uncertain{ou} \uncertain{des}} \struck{\ill{}}
\uncertain{plus} restrictif sur $\mathcal{A}$ et $\mathcal{B}$, on trouve des
catégories correspondantes d'opérations (ou constructions) sur invariants plus petites
--- \struck{\ill{}} en \uncertain{même} temps qu'on trouve des cat. « universelles » plus
petites \struck{\ill{}} \add{\uncertain{types}} invariants additifs covariants ou
contrav. possibles. \struck{\ill{}} De façon précise, il doit résulter des principes
généraux (que je n'ai pas \uncertain{détaillés}) que si \struck{\ill{}}
$F_{0} : \text{\struck{\ill{}}}\ \widehat{\Delta}^{\circ} \to \mathcal{B}_{0}$
est l'objet de $((\mathcal{A}b_{*})^{\circ})_{*} \simeq (\mathcal{A}b^{*}{}_{*})^{\circ}$
qui exprime l'objet \add{$\mathcal{B}_{0}$} semi-simplicial additif universel, alors la
\struck{\ill{}} catégorie cherchée d'opérations [\uncertain{correspondant} à un
\uncertain{type} $L$ de cat. d'indices par rapport auxquelles on se propose de
\struck{\ill{}} \uncertain{prendre} des $\varprojlim$ --- dans des cat. $\mathcal{B}$
où on regarde des inv. de nature contravariante --- resp. des $\varinjlim$ --- dans le
cas de cat. $\mathcal{A}$ où on regarde les inv. de nature cov.] est la sous-catégorie
pleine \add{$U_{L}$} de $\mathcal{B}_{0} = (\mathcal{A}b_{*})^{\circ}$ engendrée par
$F_{0}(\Delta^{n})$ grâce aux $\varprojlim$ de type $L$, \struck{et} \add{\uncertain{correspondant}}
donc à la sous-catégorie pleine de $\mathcal{A}b_{*}$ engendrée par les
$\mathbb{C}_{0}(\Delta^{n})$ \add{grâce aux} $\varinjlim$ de type $L$.
\note{« $U_{L}$ » est écrit dans la marge gauche, avec un renvoi vers « de $\mathcal{B}_{0}$ »}

\page{316}
\note{folio de l'auteur 16 ; numéro au crayon : 318}
(c'est ce qu'on vérifie directement dans certains cas remarquables, p.ex. celui où
$L = \emptyset$ et où on trouve la \add{\uncertain{sous}}-catégorie pleine de
$\mathcal{A}b_{*}$ engendrée par \add{les} $\mathbb{C}_{0}(\Delta^{n})$ grâce aux seules
sommes finies.

c) On a oublié de noter que \struck{\ill{}} dans le cas particulier
$\mathcal{B} = \mathcal{A}b$, on trouve les \uncertain{deux} types d'objets (savoir
ceux de $\mathcal{A}b_{*}$ ou de $\mathcal{A}b.$ au choix) pour exprimer les objets
\emph{inv.} contravariants \add{$F$} possibles à valeurs dans $\mathcal{A}b$ (commutant
aux $\varprojlim$) et les \emph{op.} \add{$\Omega$} qu'on peut faire sur des inv.
contrav. \struck{quelconques} \add{\uncertain{de}} \uncertain{ce} \uncertain{type},
dans une cat. \struck{\ill{}} \add{additive} à $\varprojlim$, on a déjà dit comment
\[
  \begin{aligned}
    &F \text{ correspondant à } K_{*} \in \mathcal{A}b_{*} \text{ ou } K. \in \mathcal{A}b. \\
    &\Omega \text{ \struck{\ill{}} \quad\quad\quad\quad\quad } L_{*} \in \mathcal{A}b_{*} \text{ ou } L. \in \mathcal{A}b.
  \end{aligned}
\]
comment donc s'exprime directement la correspondance entre $F$ et $\Omega$, si
$K = L$ ? On trouve
\note{dans la seconde ligne, un trait horizontal tient lieu de « correspondant à »}
\marginal{On veut ici les \uncertain{applications} \uncertain{d'objets} \ill{} \ill{} \ill{} $\mathcal{A}b_{*}$ \ill{} $\mathcal{A}b$ --- attention, dans cette formule, \uncertain{renversant} la variance par rapport à la situation --- il s'agit bien des $F$ \emph{contravariants} \ill{} à val. dans $\mathcal{A}b$ \ill{} \ill{} et non dans $(\mathcal{A}b)^{\circ}$…)}
\note{note marginale écrite en oblique entre les deux colonnes, reliée par un trait au « $(\mathcal{A}b)$ » et à la formule encadrée ; lecture très partielle}
si $F_{0} \in ((\mathcal{A}b_{*})^{\circ})_{*} = \mathcal{B}_{0*}$ \add{$\simeq (\mathcal{A}b_{*}{}^{*})^{\circ}$}
est le \struck{\ill{}} \add{\uncertain{invariant}} universel considéré, à valeurs dans
$\mathcal{B}_{0} = (\mathcal{A}b_{*})^{\circ}$, alors pour tout
$X_{*} \in \mathcal{S}s = \widehat{\Delta}$
\[
  \boxed{F(X_{*}) \simeq \Omega^{*}_{*}(\mathbb{C}_{*}(X_{*})^{\circ}) \simeq \Omega^{\cdot}.(\mathbb{C}.(X_{*})^{\circ})}
\]
où l'expression $\mathbb{C}_{*}(X_{*})$ ou $\mathbb{C}.(X_{*})$ signifie la \uncertain{même},
mais considérée comme objet non de $\mathcal{A}b_{*}$ \add{resp. $\mathcal{A}b.$}, mais de
\struck{$(\mathcal{A}b_{*})^{\circ}$}
$\mathcal{B}'^{*}$ \add{resp. $\mathcal{B}'^{\cdot}$}, où $\mathcal{B}' = \mathcal{A}b^{\circ}$,
de façon que $\Omega^{*}$ \struck{resp. $\Omega^{\cdot}$ se définisse sur $\mathcal{A}b_{*}$ resp.}
\uncertain{se} définisse sur $\mathcal{B}'^{*}$ resp. $\mathcal{B}'^{\cdot}$
($\Omega^{*}_{\mathcal{B}} : \mathcal{B}'^{*} \to \mathcal{B}'$,
$\Omega^{\cdot}_{\mathcal{B}} : \mathcal{B}'^{\cdot} \to \mathcal{B}'$) a un sens
sur $\mathbb{C}_{*}(X_{*})^{\circ}$ \add{resp. $\mathbb{C}.(X_{*})$}, et c'est
\uncertain{alors} dans $\mathcal{B} = (\mathcal{A}b)^{\circ}$ qu'on a $F(X_{*})$
(à isom. can. près).

Tous ces développements sont encore valables pour une catégorie $\Delta$ (ess. petite)
quelconque --- et aussi dans le contexte non additif --- \struck{dans \ill{}} \uncertain{avec}
les $\varinjlim$ et non les $\varprojlim$, qui, eux, \uncertain{viendraient} des
catégories de valeurs additives et la \uncertain{\ill{}} spécial de Dold--Puppe.

En fait (revenant aux types $L$ de cat. d'indices envisagés dans b)) on s'intéressera
surtout aux inv. contravariants en $X_{*}$ provenant d'objets $K_{*} \in \mathcal{A}b_{*}$
qui sont des $U_{L}$ --- i.e. qui peuvent se déduire des \uncertain{composantes}
$\mathbb{C}^{[n]}{}_{*} \in (\mathcal{A}b)_{*}$ (et des opérations semi-simpliciales
entre \add{\uncertain{icelles}}) : à l'aide des \add{$L$} opérations $\Omega_{*}$
\add{\uncertain{prenant} \uncertain{par} \uncertain{dessus}, sur des $\mathbb{C}^{n}{}_{*} \in (\mathcal{A}b_{*})^{\circ}$}
\struck{\ill{}}
\add{\uncertain{et} \uncertain{des} op. différentielles \ill{} \ill{} de telles opérations
\ill{} \ill{}} dire que l'on s'intéresse aux foncteurs \struck{\ill{}}
\note{fin de page très surchargée ; plusieurs insertions superposées, l'ordre de lecture est incertain}

\page{317}
\note{folio de l'auteur 17 ; numéro au crayon : 319}
\[
  F : X_{*} \longmapsto F(X_{*}) : \widehat{\Delta}^{\circ} \to \mathcal{A}b
\]
qui peuvent se déduire du foncteur universel
$\widehat{\Delta}^{\circ} \xrightarrow{\ F_{0}\ } (\mathcal{A}b_{*})^{\circ}$
(défini par $(F_{0})^{\circ}(X_{*}) \simeq \mathbb{C}_{*}(X_{*})$) en lui appliquant
d'opération $\Omega$ \add{\uncertain{pour} $\Omega^{*} \in U^{L}$}, ou encore qui se
déduisent du foncteur cov. universel
\[
  \widehat{\Delta} \xrightarrow{\ F_{0}{}^{\circ} = (X_{*} \mapsto \mathbb{C}_{*}(X_{*}))\ } (\mathcal{A}b_{*})
\]
\struck{qui} on \uncertain{dérivée} par $\mathbb{C}_{*}$, en lui appliquant \add{de
telles} $\Omega_{*}$. Ce \struck{\ill{}} \add{qui} se formule très
\marginal{? \quad Sous réserve que $L$ soit \uncertain{choisi} pour que les \uncertain{deux} \uncertain{formulations} \uncertain{soient} \uncertain{équivalentes} \uncertain{en} \uncertain{vertu} des \uncertain{projecteurs} !}
équivalemment en termes de $\Omega^{\cdot}$, $\Omega.$ opérant sur $(\mathcal{A}b.)^{\cdot}$
($\Omega^{\cdot} : \mathcal{A}b^{\cdot} \to \mathcal{A}b$, $\Omega. : \mathcal{A}b. \to \mathcal{A}b$),
en remplaçant $\mathbb{C}_{*}(X_{*})$ par $\mathbb{C}.(X_{*})$.

e) Quel est le rôle dans tout ceci de DP ? \struck{\ill{}} D'un côté on travaille que
l'avantage de $\mathcal{A}.$ sur $\mathcal{A}_{*}$, c'est que la description d'un objet
y est nettement plus simple --- on se perd facilement dans la description des op.
semi-simpliciales ! --- et \uncertain{donc} \uncertain{dans} \uncertain{nos} $L.$, $L_{*}$,
$L.$ est aussi nettement « moins gros » que $L_{*}$ ($\mathcal{A}b$ un \uncertain{type}
\ill{} \ill{} \uncertain{comme} \uncertain{facteur} \uncertain{direct}), $L_{*}$
\uncertain{apparaît} comme une sorte de version pléthorique de $L.$ !
\note{d) ne figure pas sur ces pages : la liste passe de c), p. 316, à e)}

Néanmoins, dès que (dans le cas de $\mathcal{A} = \mathcal{A}b_{k}$ disons)
$\mathcal{A}b_{k*}$ et $\mathcal{A}b_{k}.$ ont chacun leur structure multiplicative ---
qui ne se correspondent \emph{pas} par DP et ND --- celle de $\mathcal{A}b_{k*}$ est
plus simple à beaucoup d'égards, et s'introduit d'ailleurs de bien des façons par la
suite de façon \uncertain{absolument} impérieuse, alors que celle de $\mathcal{A}b_{k}.$
ne s'introduit pratiquement pas.
\note{le paragraphe « Néanmoins … pas » est marqué d'un trait vertical dans la marge}
De plus, pour les opérations tensorielles du type $\overset{i}{\wedge}$,
$\overset{i}{\otimes}$, $\mathrm{Sym}^{i}$, elles \uncertain{manquent} purement et simplement
dans $\mathcal{A}b_{k}.$ (sauf de les définir par transposition de structure via DP !)
\uncertain{alors} \uncertain{qu'elles} \uncertain{sont} \uncertain{évidentes} pour
sur $\mathcal{A}b_{k*}$ ! Or ces structures également joueront un rôle essentiel par
la suite.

\section{Gribouillis vrac}
\note{titre pris sur une chemise bleue (numéro au crayon 320) écrite de sa main au crayon : « \struck{Pré} Gribouillis vrac » et, plus bas, « Au Fil des Jours… » ; lecture de « Pré » douteuse. Les deux feuillets qui suivent sont des brouillons sur papier réglé, sans lien visible avec la suite paginée 1--17 qui précède}

\page{319}
\note{numéro au crayon : 321 ; feuillet de calculs en désordre, sur les puissances divisées $\Gamma^{p}$, $\Gamma^{q}$, $\Gamma^{pq}$ ; on n'en donne que les formules principales}
\[
  \text{\struck{$F_{A}(\Gamma^{pq}(X)) \longrightarrow$}} \qquad
  \Gamma^{pq}(F_{A}(X)) \xrightarrow{\ \sigma^{X}_{pq}\ } F_{A}(\Gamma^{pq}(X)) ,
  \qquad f \longmapsto \Gamma^{pq}(f) \circ u_{pq}
\]
\[
  (f^{(q)})^{(p)} \quad \Gamma^{p}(\Gamma^{q}(F_{A}(X))) \xrightarrow{\ \Gamma^{p}(\sigma^{X}_{q})\ } \Gamma^{p}F_{A}(\Gamma^{q}(X))
  \longrightarrow F_{A}(\Gamma^{p}(\Gamma^{q}(X)))
\]
\[
  (\Gamma^{q}(f) \circ u_{q})^{(p)} \longmapsto \Gamma^{p}(\Gamma^{q}(f) \circ u_{q}) \circ u_{p}
\]
\note{à droite, dans une boucle reliée par des flèches : $(\gamma^{A}_{pq} \circ \Gamma^{pq}(f) \circ u_{pq})$ et $F_{A}(\gamma^{X}_{p,q})$ ; une ligne surchargée et biffée au milieu}

\begin{tikzcd}
  A \arrow[r, "u_{pq}"] \arrow[d, "u_{p}"'] & \Gamma^{pq}(A) \arrow[r, "\Gamma^{pq}(f)"] & \Gamma^{pq}(X) \arrow[d, "\gamma^{X}_{pq}"] \\
  \Gamma^{p}(A) \arrow[r, "\Gamma^{p}(u_{q})"'] & \Gamma^{p}(\Gamma^{q}(A)) \arrow[r, "{\Gamma^{p}(\Gamma^{q}(f))}"'] & \Gamma^{p}(\Gamma^{q}(X))
\end{tikzcd}

\note{deux flèches obliques, partant de $\Gamma^{p}(A)$ et de $\Gamma^{p}(\Gamma^{q}(A))$ vers la ligne du haut, sont à demi biffées}
\[
  f : A \to X, \qquad \Gamma^{q}f : \Gamma^{q}A \to \Gamma^{q}X, \qquad \Gamma^{q}f \circ u_{q} \ (\text{de } A)
\]
$X = A$, $f = \mathrm{id}_{A}$ :

\begin{tikzcd}
  A \arrow[r, "u_{pq}"] \arrow[d, "u_{p}"'] & \Gamma^{pq}(A) \arrow[d, "\gamma^{A}_{pq}"] \\
  \Gamma^{p}(A) \arrow[r, "\Gamma^{p}(u_{q})"'] & \Gamma^{p}(\Gamma^{q}(A))
\end{tikzcd}

\note{ce carré est entouré d'un trait}

\begin{tikzcd}
  A \otimes A \arrow[r, "u_{p} \otimes u_{q}"] & \Gamma^{p}(A) \otimes \Gamma^{q}(A) \arrow[d, "\mu^{A}_{p,q}"] \\
  A \arrow[u, "\Delta"] \arrow[r, "c_{pq}\, u_{p+q}"'] & \Gamma^{p+q}(A)
\end{tikzcd}

\note{le nom de la flèche verticale de droite est surchargé ; $\mu$ est une lecture douteuse, de même que « $c_{pq}$ » sous la flèche du bas}
\[
  \boxed{u_{p} \cdot u_{q} = c_{pq}\, u_{p+q}} \quad \text{dans} \quad
  \mathrm{Hom}_{1}(A, \Gamma^{\cdot} A) \simeq (\mathrm{Hom}_{1}(\mathrm{Sym}^{\cdot}\check{A}, \check{A}))
  \qquad \boxed{\begin{array}{l} u_{0} = 1 \\ u_{1} = \mathrm{id} \end{array}}
\]
\uncertain{qui} \uncertain{est} une alg. \uncertain{grâce} \uncertain{à} la structure de cog. sur $A$ et la structure d'alg. de $\Gamma^{\cdot} A$
\note{au-dessus de $\mathrm{Hom}_{1}$, une flèche venant de « $k\{T\}$ »}

\page{320}
\note{numéro au crayon : 322 ; la moitié supérieure de la page est barrée d'un long trait oblique}
\[
  X \xrightarrow{\ \text{dérivé}\ } Y, \qquad \Gamma^{n}(X) \nearrow, \qquad F(X) \overset{?}{\longrightarrow} F(Y)
\]
\[
  F(\Gamma^{n}X) \longrightarrow F(Y), \qquad \Gamma^{n}(F(X)) \longrightarrow F(\Gamma^{n}X)
\]
\note{ce premier bloc est séparé du suivant par un trait horizontal}
\[
  F_{A}(X) = \mathrm{Hom}_{1}(A, X)
\]
\[
  \Gamma^{n}F_{A}(X) \overset{?}{\longrightarrow} F_{A}(\Gamma^{n}X)
\]
\[
  \Gamma^{n}\mathrm{Hom}_{1}(A, X) \xrightarrow{\ \beta\,?\ } \mathrm{Hom}_{1}(A, \Gamma^{n}X)
\]
\struck{$\beta(u^{(n)}) \to \ldots\ \Gamma^{n}(u)$} fonctorielle en \ill{}
\[
  X \longrightarrow X \otimes X
\]
\[
  F_{C}(X \otimes Y) \xleftarrow{\ f\ } F_{A}(X) \times F_{B}(Y) \qquad \text{bilin., fonction en } X, Y
\]
si $X = Y = A$, c'est \ill{} $\mathrm{id}_{A}$, $\mathrm{id}_{A}$ donne une \uncertain{flèche}
\[
  A \longrightarrow A \otimes A
\]
\begin{enumerate}
  \item[a)] \uncertain{associativité} \add{coassoc.} : \uncertain{morphisme} de la \uncertain{cohom.} \ill{}
  \item[b)] \uncertain{commutative} \add{\uncertain{cocommutative}}
  \item[c)] compatible avec $1$ \quad i.e. $k \to F_{A}(\mathbb{1}) = \mathrm{Hom}(A, \mathbb{1})$
\end{enumerate}
\uncertain{i.e.} \uncertain{counité} \quad $A \xrightarrow{\ n\ } \mathbb{1}$
\note{l'étiquette de la flèche $X \to Y$ est lue « dérivé » avec doute ; celle de $A \to \mathbb{1}$ est lue « $n$ »}

\end{document}
